% !TeX encoding = UTF-8
% !TeX program = pdflatex
% !BIB program = biber
%
% LNI (Lecture Notes in Informatics, Gesellschaft für Informatik) artifact paper.
% Build requires the GI LNI document class (lni.cls) and biber; see README.md.
% Modelled on the LNI setup used in publications/pub_rse_methodology (_quarto-lni.yml).
%
\documentclass[english,biblatex]{lni}

\usepackage{graphicx}
\usepackage{booktabs}
\usepackage{url}

\addbibresource{artifact.bib}

\begin{document}

\title[RSE in the DELFI Community --- Artifact]{%
Artifact: A Reproducible, Database-Free Pipeline for Detecting Research-Software
Paradigms in the DELFI Community}

% LNI author interface (authblk-style): \author[affil#]{name}{email}{orcid}
\author[1]{Julian Dehne}{dehne3@uni-potsdam.de}{0000-0001-9265-9619}
\author[2]{Fabian Lochner}{fabian.lochner@example.org}{}
\author[1]{Axel Wiepke}{axel.wiepke@uni-potsdam.de}{}
\affil[1]{University of Potsdam, An der Bahn 2, 14476 Potsdam, Germany}
\affil[2]{Affiliation of the second author}

\maketitle

\begin{abstract}
This artifact accompanies the DELFI~2026 paper \emph{Research Software Engineering
in the Learning Technologies: a replication study of research paradigms in the
DELFI community}. It provides a self-contained, file-based pipeline that
reproduces the study's label aggregation, inter-coder reliability (ICR) metrics,
descriptive figures and secondary bias analysis directly from the shipped
intermediate data --- \emph{without} access to the copyrighted DELFI paper PDFs
or to a commercial LLM API. The six pipeline steps are orchestrated by a single
management script; the two steps that require the paper texts (preprocessing,
human re-annotation) and the one that requires an LLM API token (the experiments)
are skipped by default. A programmatic guard enforces that no paper full text is
ever contained in the released data. We describe the artifact's structure, how to
run it, the reproducibility guarantees, and the constraints that motivated its
design for the artifact-evaluation track.
\end{abstract}

\begin{keywords}
Research Software Engineering \and replication \and large language models \and
inter-coder reliability \and reproducibility \and DELFI
\end{keywords}

\section{Introduction}
The companion study~\cite{dehne2026rseLearningTechnologies} replicates and extends
a 2017 analysis of research-software paradigms across the DELFI proceedings, using
three open-weight large language models (LLMs) to annotate $n=798$ papers and
validating the result against the earlier human-coded reference study. This artifact packages the \emph{evidence and
the computation} behind those results so that a reviewer can re-run the analysis
end-to-end on a laptop in minutes.

The central design constraint is legal: the analyzed DELFI papers are copyrighted
and may not be redistributed. The artifact therefore ships only derived data
(LLM labels, justifications, human re-annotations and bibliographic metadata) and
is explicitly built so that the two paper-dependent steps and the API-dependent
step can be skipped while still reproducing every published number and figure.

\section{Artifact overview}
The artifact is a Python package (\texttt{pipeline/}) plus shipped intermediate
data (\texttt{data/}) and this paper (\texttt{artifact-paper/}). It is database-free:
whereas the original study stored papers in MySQL, every step here reads and writes
CSV files under \texttt{data/}. We claim the \emph{Available} and \emph{Reusable}
properties: the package is permanently archived, documented, and parameterised so
that it can also be re-run on new inputs.

\section{Repository structure}
\begin{verbatim}
rse_in_delfi_artifact/
  run_pipeline.py         # management script (skip-aware)
  pipeline/               # config + the six steps + full-text guard
  data/                   # shipped intermediate products (no full texts)
  artifact-paper/         # this LNI document
\end{verbatim}

\section{The pipeline}
The pipeline has six steps (Table~\ref{tab:steps}). Steps~3, 5 and~6 run from the
shipped data and constitute the default reviewer run; steps~1, 2 and~4 are
\emph{skipped by default} because they need resources reviewers do not have.

\begin{table}[h]
\centering
\caption{Pipeline steps and default execution policy.}
\label{tab:steps}
\begin{tabular}{clll}
\toprule
\# & Step & Default & Requires \\
\midrule
1 & Preprocessing (PDF $\rightarrow$ text)      & \textbf{skipped} & paper PDF folder \\
2 & LLM experiments                              & \textbf{skipped} & SAIA/KISSKI API token \\
3 & Descriptive analysis (figures)               & runs             & shipped data \\
4 & Human re-annotation                          & \textbf{skipped} & paper PDF folder (interactive) \\
5 & Label aggregation (intra + inter, ICR)       & runs             & shipped data \\
6 & Secondary analysis (paradigm biases)         & runs             & shipped data \\
\bottomrule
\end{tabular}
\end{table}

\paragraph{Data flow.} Step~1 extracts text from the paper PDFs; step~2 sends each
paper to three LLMs (three runs each) via the KISSKI SAIA API and stores the
labels and justifications. Step~5a aggregates the three runs per model
(majority vote) and computes intra-model ICR; step~4 elicits a human label for the
papers on which all three models disagree; step~5b aggregates across models
(majority vote, ties broken by the human label), reports cross-model ICR and
validates against the 2017 reference study. Steps~3 and~6 produce the figures and
the secondary bias analysis from the final aggregated labels.

\paragraph{Management script.} \texttt{run\_pipeline.py} runs the default steps,
accepts \texttt{-{}-paper-folder} as a parameter to enable steps~1 and~4, and
\texttt{-{}-run-experiments} (with a SAIA token in the environment) to enable
step~2. A no-full-text guard runs before and after every invocation.

\section{Reproducing the results}
\begin{verbatim}
python -m pip install -r requirements.txt
python run_pipeline.py            # steps 3, 5, 6 from shipped data
\end{verbatim}
This regenerates \texttt{data/inter\_LLM/df\_llm\_experiments\_final\_aggregated\_*.csv}
($798\times185$), the ICR tables, the paradigm-distribution and label-trend figures,
and the combined-label bias tables.

\section{Reproduced results}
Table~\ref{tab:icr} shows the cross-model ICR recomputed by step~5 on the shipped
data, reported as Krippendorff's $\alpha$~\cite{krippendorff2004reliability} (and,
for binary labels, complemented by Gwet's AC1~\cite{gwet2008ac1} in the released
tables). All three models disagreed completely on 0, 6, 17 and 9 papers for the
process, outcome, design and acceptance paradigms respectively; those cases were
resolved with the shipped human labels.

\begin{table}[h]
\centering
\caption{Cross-model inter-coder reliability (Krippendorff's $\alpha$), recomputed.}
\label{tab:icr}
\begin{tabular}{lc}
\toprule
Label group & Krippendorff's $\alpha$ \\
\midrule
Binary labels (average)  & 0.51 \\
Ordinal labels (average) & 0.66 \\
All labels (average)     & 0.60 \\
\midrule
Reference-study validation (binary, average) & 0.33 \\
\bottomrule
\end{tabular}
\end{table}

\section{Data and the no-full-text guarantee}
The released data contain only labels, LLM/human justifications and bibliographic
metadata. The original annotation script dropped the \texttt{abstract},
\texttt{text} and \texttt{references} columns before persisting results ``for legal
reasons''; the artifact preserves this invariant and enforces it programmatically.
\texttt{pipeline/fulltext\_guard.py} scans every shipped table and fails if a banned
full-text column is present or if any cell exceeds a length threshold incompatible
with a short justification. Step~1's full-text output is written only to a
\texttt{.gitignore}'d directory.

\section{Requirements and environment}
Python~3.11+, \texttt{pandas}, \texttt{numpy}, \texttt{scikit-learn},
\texttt{krippendorff}, \texttt{statsmodels}, \texttt{matplotlib} and
\texttt{openpyxl}; see \texttt{requirements.txt}. \texttt{pymupdf} is needed only to
re-run step~1, and an OpenAI-compatible client plus a SAIA token only to re-run
step~2.

\section{Limitations}
Steps~1 and~2 cannot be exercised by reviewers without the copyrighted PDFs and an
API token; their outputs are shipped and verified downstream instead. LLM
annotations are sampled at temperature~0 with a fixed seed, but exact
reproducibility of step~2 ultimately depends on the provider's model snapshots.

\nocite{rseInDelfiArtifact}
\printbibliography

\end{document}
